\DocumentMetadata{
  pdfversion=2.0,
  pdfstandard=ua-2,
  lang=en-US,
  tagging=on,
  tagging-setup={math/setup=mathml-SE,tabsorder=structure}
}
\documentclass[11pt,letterpaper]{article}
\usepackage[margin=0.68in,includefoot]{geometry}
\usepackage{newcomputermodern}
\usepackage{amsmath,amscd,mathtools}
\usepackage{tikz,adjustbox}
\providecommand{\square}{\mdlgwhtsquare}
\usepackage[hidelinks]{hyperref}
\hypersetup{
  pdftitle={Algebra PhD Exam, May 2013},
  pdfauthor={University of Florida Department of Mathematics},
  pdfsubject={Graduate mathematics examination},
  pdfdisplaydoctitle=true,
  bookmarksopen=true
}
\setcounter{secnumdepth}{0}
\setlength{\parindent}{0pt}
\setlength{\parskip}{0.28em}
\setlength{\emergencystretch}{3em}
\setlength{\leftmargini}{2.15em}
\setlength{\leftmarginii}{2.65em}
\setlength{\leftmarginiii}{3em}
\pagestyle{empty}
\raggedbottom
\allowdisplaybreaks
\NewTaggingSocketPlug{block/list/label}{examproblem}{%
  \tagstructbegin{tag=Lbl}\tagstructend%
  \tagstructbegin{tag=\LBody}%
  \tagstructbegin{tag=\ProblemHeadingTag,title-o={\ProblemHeadingText}}%
  \tagmcbegin{tag=\ProblemHeadingTag,actualtext={\ProblemHeadingText}}%
  #2%
  \tagmcend%
  \tagstructend%
}
\newenvironment{examproblems}[2]{%
  \def\ProblemHeadingTag{#2}%
  \AssignTaggingSocketPlug{block/list/label}{examproblem}%
  \begin{enumerate}%
  \setcounter{enumi}{#1}%
  \renewcommand{\labelenumi}{\textbf{\theenumi.}}%
  \setlength{\topsep}{0.35em}%
  \setlength{\partopsep}{0pt}%
  \setlength{\itemsep}{0.48em}%
  \setlength{\parsep}{0.18em}%
}{\end{enumerate}}
\newenvironment{examsubparts}{%
  \AssignTaggingSocketPlug{block/list/label}{default}%
  \begin{enumerate}%
  \setlength{\topsep}{0.18em}%
  \setlength{\partopsep}{0pt}%
  \setlength{\itemsep}{0.16em}%
  \setlength{\parsep}{0.1em}%
}{\end{enumerate}}
\newenvironment{examinstructions}{%
  \par\begingroup\itshape
}{%
  \par\endgroup\vspace{0.18em}
}
\makeatletter
\renewcommand\subsection{\@startsection{subsection}{2}{0pt}%
  {-1.25ex plus -.25ex minus -.1ex}%
  {0.55ex plus .1ex}%
  {\normalfont\large\bfseries}}
\renewcommand{\@maketitle}{%
  \newpage\null\vskip -1em%
  {\footnotesize\bfseries
    \href{https://www.ufl.edu/}{University of Florida}, \href{https://math.ufl.edu/}{Department of Mathematics}\par}%
  \vskip -0.5em%
  \tagstructbegin{tag=H1,title={Algebra PhD Exam, May 2013}}%
  \tagpdfparaOff%
  \tagmcbegin{tag=H1}%
    {\LARGE\bfseries\href{https://gma.math.ufl.edu/exams/algebra-phd/}{Algebra PhD Exam}, May 2013\par}%
  \tagmcend%
  \tagpdfparaOn%
  \tagstructend%
  \vskip 0.25em%
  \tagmcbegin{artifact}%
    \hrule height 1pt%
  \tagmcend%
  \vskip 1em%
}
\makeatother
\title{Algebra PhD Exam, May 2013}
\author{}
\date{}
\begin{document}
\maketitle
\thispagestyle{empty}
\begin{examinstructions}
Answer seven problems. Write your answers clearly in complete English sentences. You may quote results (within reason) as long as you state them clearly.
\end{examinstructions}
\begin{examproblems}{0}{H2}
\def\ProblemHeadingText{Problem 1.}
\item
Let~\(K\) be a field and let \(f(x) \in K[x]\) be a non-zero polynomial.
\begin{examsubparts}
\item[\textnormal{(a)}]
Define what is a splitting field for \(f(x)\) over~\(K\).
\item[\textnormal{(b)}]
If \(F_1\) and \(F_2\) are splitting fields for \(f(x)\) over~\(K\) prove, from your definition, that \(F_1\) and \(F_2\) are isomorphic fields.
\end{examsubparts}
\def\ProblemHeadingText{Problem 2.}
\item
Let \(F=\mathbb{Q}(\sqrt{3},\sqrt{5})\) be the smallest subfield of~\(\mathbb{C}\) which contains \(\sqrt{3}\) and \(\sqrt{5}\). How many subfields does~\(F\) have?
\def\ProblemHeadingText{Problem 3.}
\item
Let~\(A\) and~\(B\) be finite abelian groups and suppose that \(|A|\) and \(|B|\) are relatively prime. Let \(T=A\otimes_{\mathbb{Z}}B\). Show that \(|T|=1\).
\def\ProblemHeadingText{Problem 4.}
\item
Let~\(C\) be a category and let \((A_i)_{i\in I}\) be a family of objects of~\(C\).
\begin{examsubparts}
\item[\textnormal{(a)}]
Define what is a coproduct of \((A_i)_{i\in I}\) in~\(C\).
\item[\textnormal{(b)}]
Prove that any two coproducts of \((A_i)_{i\in I}\) in~\(C\) are equivalent (some texts would say isomorphic) in~\(C\).
\end{examsubparts}
\def\ProblemHeadingText{Problem 5.}
\item
Let~\(R\) be a ring.
\begin{examsubparts}
\item[\textnormal{(a)}]
Define what it means for an~\(R\)-module~\(M\) to be projective.
\item[\textnormal{(b)}]
Prove that if
\[
0\to A\to B\to P\to 0
\]
 is an exact sequence of~\(R\)-modules and~\(P\) is projective then the exact sequence splits.
\end{examsubparts}
\def\ProblemHeadingText{Problem 6.}
\item
Let~\(F\) be a free group with free generators~\(a\) and~\(b\). Let~\(N\) be the smallest normal subgroup of~\(F\) such that \(a^2\in N\) and \(b^2\in N\). Prove that \(F/N\) is infinite.
\def\ProblemHeadingText{Problem 7.}
\item
State and prove the Lying-Over Theorem on prime ideals and ring extensions.
\def\ProblemHeadingText{Problem 8.}
\item
Prove the following version of Nakayama’s Lemma. Let~\(A\) be a commutative ring with identity, let~\(M\) be a finitely generated~\(A\)-module, and let~\(I\) be an ideal contained in the intersection of all the maximal ideals of~\(A\). Show that if \(IM=M\) then \(M=\{0\}\).
\def\ProblemHeadingText{Problem 9.}
\item
An invertible ideal in an integral domain that is a local ring is principal.
\def\ProblemHeadingText{Problem 10.}
\item
Let \(\overline{\mathbb{Q}}\) be an algebraic closure of~\(\mathbb{Q}\) and \(v\in\overline{\mathbb{Q}}\) with \(v\notin\mathbb{Q}\). Let~\(E\) be a subfield of \(\overline{\mathbb{Q}}\) maximal with respect to the condition that \(v\notin E\). Prove that every finite-dimensional extension of~\(E\) is cyclic.
\def\ProblemHeadingText{Problem 11.}
\item
Classify (up to ring isomorphism) all semisimple rings of order \(720\).
\end{examproblems}
\end{document}
